% !TEX program = xelatex
% ===========================================================================
%  SCX协议社区宪法：三层治理架构、维护者轮换博弈论与g=0宣言协议
%  SCX Protocol Community Constitution: Three-Layer Governance,
%  Maintainer Rotation Game Theory, and the g=0 Declaration Protocol
%  Date: 2026-07-02
%  Classification: SCX Internal — Public Release Approved
% ===========================================================================

\documentclass[12pt,a4paper]{ctexart}

% ---- Chinese & Font Support ----
\usepackage[UTF8]{ctex}
\setCJKmainfont{SimSun}[BoldFont=SimHei,ItalicFont=KaiTi]
\setmainfont{Times New Roman}

% ---- Mathematics ----
\usepackage{amsmath,amssymb,amsthm,mathtools}
\usepackage{bm}
\usepackage{bbm}
\usepackage{dsfont}

% ---- Theorem Environments ----
\newtheorem{theorem}{定理}[section]
\newtheorem{proposition}[theorem]{命题}
\newtheorem{corollary}[theorem]{推论}
\newtheorem{lemma}[theorem]{引理}
\newtheorem{definition}[theorem]{定义}
\theoremstyle{remark}
\newtheorem{remark}[theorem]{注记}
\newtheorem{example}[theorem]{示例}

% ---- Layout ----
\usepackage[margin=2.5cm]{geometry}
\usepackage{booktabs}
\usepackage{graphicx}
\usepackage{hyperref}
\usepackage{enumitem}
\usepackage[capitalise,noabbrev]{cleveref}
\usepackage{xcolor}
\usepackage{tikz}
\usepackage{float}
\usepackage{longtable}
\usepackage{makecell}
\usepackage{multirow}
\usepackage{setspace}
\usepackage{physics}
\usepackage{fancyhdr}
\usepackage{tcolorbox}
\usepackage{array}

\onehalfspacing

% ---- Hyperref ----
\hypersetup{
    colorlinks=true,
    linkcolor=blue,
    citecolor=blue,
    urlcolor=blue,
    pdftitle={SCX协议社区宪法：三层治理架构、维护者轮换博弈论与g=0宣言协议},
    pdfauthor={SCX},
    pdfsubject={SCX Protocol Community Constitution},
}

% ---- Custom Commands ----
\newcommand{\SCX}{\mathsf{SCX}}
\newcommand{\Yajie}{\mathsf{Yajie}}
\newcommand{\Spring}{\mathsf{Spring}}
\newcommand{\Situs}{\mathsf{Situs}}
\newcommand{\Cercis}{\mathsf{Cercis}}
\newcommand{\R}{\mathbb{R}}
\newcommand{\E}{\mathbb{E}}
\newcommand{\Pbb}{\mathbb{P}}
\newcommand{\Var}{\mathrm{Var}}
\newcommand{\indicator}{\mathbf{1}}
\newcommand{\indep}{\perp\!\!\!\perp}

% Game theory notation
\newcommand{\payoff}[1]{u_{#1}}
\newcommand{\deviate}{\delta}
\newcommand{\auditCost}{c_{\text{audit}}}
\newcommand{\rotationPeriod}{T}
\newcommand{\numMaintainers}{M}
\newcommand{\biasParam}{g}
\newcommand{\kernelSet}{\mathcal{K}}
\newcommand{\contributorSet}{\mathcal{C}}
\newcommand{\observerSet}{\mathcal{O}}

% ---- Color Definitions ----
\definecolor{scxblue}{RGB}{0,51,102}
\definecolor{scxred}{RGB}{180,20,20}
\definecolor{scxgold}{RGB}{200,160,20}
\definecolor{scxgreen}{RGB}{20,120,20}
\definecolor{scxgray}{RGB}{100,100,100}
\definecolor{critiquered}{RGB}{200,30,30}
\definecolor{insightblue}{RGB}{30,30,180}

% ---- Honest Critique — SCX House Style ----
\newtheorem{honesthit}{诚实暴击}[section]
\renewcommand{\thehonesthit}{\textcolor{critiquered}{\textbf{[暴击 \arabic{honesthit}]}}}
\newenvironment{hitbox}{%
  \begin{center}
  \color{critiquered}
  \begin{minipage}{0.92\textwidth}
  \begin{honesthit}
}{%
  \end{honesthit}
  \end{minipage}
  \end{center}
}

% ---- Protocol Governance Note ----
\newtheorem{govnote}{治理注记}[section]
\renewcommand{\thegovnote}{\textcolor{scxblue}{\textbf{[治理透视 \arabic{govnote}]}}}
\newenvironment{govbox}{%
  \begin{center}
  \color{scxblue}
  \begin{minipage}{0.92\textwidth}
  \begin{govnote}
}{%
  \end{govnote}
  \end{minipage}
  \end{center}
}

% ---- Constitutional Principle ----
\newtheorem{principle}{宪章原则}[section]
\renewcommand{\theprinciple}{\textcolor{scxgold}{\textbf{[原则 \arabic{principle}]}}}

% ---- Title ----
\title{\textbf{SCX协议社区宪法}\\[8pt]
       \large 三层治理架构、维护者轮换博弈论与$g=0$宣言协议\\[6pt]
       \large SCX Protocol Community Constitution:\\[4pt]
       Three-Layer Governance, Maintainer Rotation Game Theory,\\and the $g=0$ Declaration Protocol\\[6pt]
       \normalsize ——社区是定理的守护者，定理不依赖任何个人\\[4pt]
       --- The Community Guards the Theorem; the Theorem Depends on No Individual}
\author{SCX 协议社区治理工作组 \\[3pt] \small SCX Protocol Community Governance Working Group}
\date{2026年7月2日 \\ July 2, 2026}

\begin{document}

\maketitle

% ===================================================================
\begin{abstract}
\noindent
This paper establishes the constitutional framework for the SCX protocol community.
We formalize a three-layer governance architecture consisting of the Kernel
($\kernelSet$, 2--5 members), Contributors ($\contributorSet$, approximately 20
active members), and Observers ($\observerSet$, unlimited). The architecture is
anchored in four mathematically provable pillars: (1) the $g=0$ declaration
protocol, which requires every maintainer to publicly renounce ownership over the
protocol; (2) $M>1$ mutual audit of every maintainer's bias parameter $g$, ensuring
that deviation is detectable; (3) finite rotation cycles that bound the probability
of undetected deviation by $e^{-2M\Delta^2}$ via Hoeffding's inequality; and
(4) a unanimous consent rule for kernel membership, preventing capture by any
coalition.

We provide a game-theoretic analysis of the maintainer rotation mechanism, proving
that the unique subgame-perfect equilibrium of the repeated rotation game is for
every maintainer to maintain $g=0$ throughout their rotation term. The conflict
resolution protocol specifies a three-stage escalation path: extended kernel
deliberation (72 hours), contributor-layer public deliberation with non-binding
advisory vote, and independent third-party $M>1$ audit. Emergency procedures
enable any contributor to trigger an emergency audit within 48 hours upon
presenting evidence of $g \neq 0$ by a kernel member.

The community is constitutionally defined as a non-profit, open-source, zero-ownership
protocol maintenance network. It is not a company, not a fan club, not a
paid-membership organization, and not any individual's personal brand. The
community belongs to the theorem — and the theorem alone.

\medskip
\noindent\textbf{Keywords:} protocol governance, community constitution, maintainer
rotation, game theory, $g=0$ declaration, three-layer architecture, mutual audit,
conflict resolution, emergency procedures, non-profit protocol maintenance

\medskip
\noindent\textbf{摘要：}
本文确立SCX协议社区的宪法框架。我们将社区形式化为三层治理架构：内核层
（$\kernelSet$，2--5人）、贡献者层（$\contributorSet$，约20名活跃成员）和观察者层
（$\observerSet$，不限人数）。该架构以四个数学上可证明的支柱为基础：(1) $g=0$宣言协议，
要求每位维护者公开放弃对协议的所有权主张；(2) 对每位维护者的偏差参数$g$进行$M>1$
互相审计，确保偏差可检测；(3) 有限轮换周期通过Hoeffding不等式将未检测偏差的概率限制在
$e^{-2M\Delta^2}$；(4) 内核成员加入的全票通过规则，防止任何联盟夺取控制权。

我们提供了维护者轮换机制的博弈论分析，证明重复轮换博弈的唯一子博弈完美均衡是每位
维护者在整个轮换任期内始终保持$g=0$。冲突解决协议规定了三级升级路径：内核内部延长
讨论（72小时）、贡献者层面公开讨论与非约束性建议投票、以及独立第三方$M>1$审计。
紧急程序允许任何贡献者在提交内核成员$g \neq 0$的证据后，于48小时内触发紧急审计。

社区在宪法上被定义为一个非盈利、开源、零所有权的协议维护网络。它不是公司，不是粉丝群，
不是付费会员制组织，不是任何个人的品牌。社区属于定理——且仅属于定理。

\bigskip
\noindent\textbf{关键词：} 协议治理，社区宪法，维护者轮换，博弈论，$g=0$宣言，
三层架构，互相审计，冲突解决，紧急程序，非盈利协议维护
\end{abstract}

\newpage
\tableofcontents
\newpage

% ===================================================================
% SECTION 1: WHY A COMMUNITY
% ===================================================================
\section{为什么需要社区：定理、代码与人的三角结构}
\section{Why a Community: The Triangle of Theorem, Code, and People}

\subsection{问题的提出}

SCX协议已经具备了三个核心要素中的两个：\textit{定理}——从SCX因果审计框架到\Yajie{}多专家纳什-帕累托均衡，
从\Spring{}配置空间收敛到\Situs{}审计不变量，完备的数学基础已经建立；
\textit{代码}——所有核心组件均已实现并开源，包括校准锚点维护和审计日志生成。
但第三个要素仍然缺失：\textit{人}。

\begin{govbox}
\textbf{历史教训：协议的死亡模式。} TCP/IP没有社区就死了。Linux没有社区就死了。
Python没有社区就死了。所有存活至今的基础设施协议的共同特征不是技术优雅，
而是有一个持续的、自我更新的维护者网络。技术可以完美——但如果维护者消失，
协议就变成一具无法进化的尸体。
\end{govbox}

SCX需要的不是用户——用户自然会来。SCX需要的不是崇拜者——崇拜者只会点头。
SCX需要的是一群\textbf{维护者}：一群能够审计定理、审查代码、复现结果、
并最终在维护者离开后接替其位置的人。

\begin{definition}[协议社区的生命力函数]
设协议$P$在时刻$t$的维护者集合为$\mathcal{M}(t)$，定义社区生命力：
\begin{equation}
    \Lambda(P, t) = \min_{m \in \mathcal{M}(t)} \left|\frac{\partial \mathcal{M}(t)}{\partial t}\right|
    \cdot \indicator{\exists \text{ succession path for each } m}
\end{equation}
协议存活当且仅当$\Lambda(P, t) > 0$对所有$t$成立——即每个维护者都有可验证的继任路径，
且维护者集合不会缩减至空。
\end{definition}

\subsection{社区不是、社区是}

在定义社区架构之前，必须首先澄清社区\textit{不是}什么。这是一个宪法边界问题——
边界的清晰程度直接决定了社区能否抵抗捕获和腐化。

\begin{table}[H]
\centering
\caption{社区边界的宪法界定 \\ Constitutional Boundary Definition of the Community}
\begin{tabular}{>{\color{critiquered}\bfseries}p{4cm}|>{\color{scxgreen}\bfseries}p{5cm}}
\toprule
\textbf{不是 / NOT} & \textbf{是 / IS} \\
\midrule
公司 (Company) & 协议维护网络 (Protocol Maintenance Network) \\
粉丝群 (Fan Club) & 审计者社区 (Auditor Community) \\
付费会员制 (Paid Membership) & 开源贡献者网络 (Open-Source Contributor Network) \\
任何个人的品牌 (Personal Brand) & 定理的共同守护者 (Collective Guardians of the Theorem) \\
\bottomrule
\end{tabular}
\end{table}

\begin{principle}
\textbf{社区边界不可侵蚀原则。} 以上四项否定界定是宪法性约束。任何将社区转化为公司、
粉丝群、付费组织或个人品牌的尝试，均构成对社区宪法的根本违反。该违反本身即构成
$g \neq 0$的证据——因为此类尝试必然以牺牲协议中立性为代价来获取个人或机构利益。
\end{principle}

% ===================================================================
% SECTION 2: THREE-LAYER ARCHITECTURE
% ===================================================================
\section{社区架构：三层治理模型}
\section{Community Architecture: The Three-Layer Governance Model}

\subsection{结构概览}

SCX社区采用三层同心圆架构。层级之间的移动取决于两个因素：
贡献的客观可验证性和偏差参数$g$的审计结果。层级不是荣誉头衔——层级是\textit{信任假设的级别}：
层级越高，信任假设越低，审计要求越严格。

\begin{table}[H]
\centering
\caption{SCX社区三层架构 \\ SCX Community Three-Layer Architecture}
\renewcommand{\arraystretch}{1.4}
\begin{tabular}{cclp{3.5cm}p{4cm}}
\toprule
\textbf{层} & \textbf{规模} & \textbf{成员} & \textbf{职责} & \textbf{门槛} \\
\textbf{Layer} & \textbf{Size} & \textbf{Who} & \textbf{Responsibilities} & \textbf{Threshold} \\
\midrule
\multirow{2}{*}{\textbf{内核 Kernel}} 
    & \multirow{2}{*}{2--5} 
    & 初始维护者
    & \makecell[l]{轮换维护\Yajie{} API校准锚点 \\ (CEC基准值)}
    & \multirow{2}{*}{\makecell[l]{$M>1$审计通过 \\ $g=0$公开声明}} \\
\cline{3-4}
    & & 资深审计者
    & 审核新维护者候选人 & \\
\midrule
\multirow{2}{*}{\textbf{贡献者 Contributors}} 
    & \multirow{2}{*}{$\sim$20} 
    & 活跃开发者
    & \makecell[l]{提交代码、审查论文 \\ 发现漏洞、复现审计结果}
    & \multirow{2}{*}{\makecell[l]{至少一次有效贡献 \\ 被内核接受}} \\
\cline{3-4}
    & & 审稿人
    & 参与 hostile review & \\
\midrule
\textbf{观察者 Observers} 
    & 不限 
    & 任何人
    & \makecell[l]{使用\Spring{}、运行\Yajie{} \\ 提issue、参与讨论}
    & 零 \\
\bottomrule
\end{tabular}
\end{table}

\subsection{内核层 $\kernelSet$}

内核是SCX协议的\textit{操作核心}——不是权力核心。内核成员不拥有协议，
不控制协议的方向，不享有任何高于其他层级成员的特权。内核成员唯一的特殊职责是：
\textbf{轮换维护\Yajie{} API的校准锚点}，即CEC（Consensus Evaluation Criterion）基准值。

\begin{definition}[内核层形式化定义]
\begin{align}
    \kernelSet &= \{k_1, \ldots, k_{\numMaintainers}\}, \quad 2 \leq \numMaintainers \leq 5 \\
    \text{条件}: \quad &\forall k_i \in \kernelSet: \\
    &(1)\; g(k_i) = 0 \text{ 经 $M>1$ 独立审计确认} \nonumber \\
    &(2)\; \text{公开 $g=0$ 声明已签名并发布} \nonumber \\
    &(3)\; \text{处于轮换任期内，任期 $T \leq T_{\max}$} \nonumber
\end{align}
其中$g(k_i)$是维护者$k_i$的偏差参数，$T_{\max}$是最大轮换任期。
\end{definition}

\subsection{贡献者层 $\contributorSet$}

贡献者是SCX社区的\textit{免疫系统}。他们通过持续的代码审查、论文 hostile review、
实验复现和漏洞发现，构成了对内核的第一道外部审计防线。贡献者不能修改CEC基准值，
但可以\textit{公开质疑}任何校准决策，并要求内核提供审计日志支持。

\begin{definition}[贡献者层准入条件]
个体$c$有资格成为贡献者，当且仅当：
\begin{equation}
    c \in \contributorSet \iff \exists \text{ 至少一次有效贡献 } \iota \text{ 被 } \kernelSet \text{ 接受}
\end{equation}
其中``有效贡献''的判定标准由内核公开维护，且该判定本身可被任何观察者审计。
\end{definition}

\subsection{观察者层 $\observerSet$}

观察者层是SCX社区的\textit{开放性}的制度保证。任何人——无论背景、资历、地域——
均可零门槛成为观察者。观察者可以使用\Spring{}、运行\Yajie{}、提交issue、
参与公开讨论。观察者层的存在使得社区永远不可能成为封闭的精英俱乐部。

\begin{govbox}
\textbf{观察者层的战略意义。} 观察者层不是被动的旁观者池——它是贡献者和内核的
\textit{人才蓄水池}。每一个内核成员都是从观察者起步、经过贡献者阶段成长起来的。
零门槛意味着社区永远不会面临``维护者后继无人''的困境——只要协议有价值，
就有人在使用；只要有人在使用，就有潜在的维护者。
\end{govbox}

% ===================================================================
% SECTION 3: THE PATH TO THE KERNEL
% ===================================================================
\section{进入内核的路径：从贡献者到守护者}
\section{The Path to the Kernel: From Contributor to Guardian}

\subsection{五阶段晋升流程}

内核成员资格不是授予的——是\textit{经过审计证明的}。晋升路径包含五个不可跳过的阶段，
每个阶段都产生可公开审计的记录。

\begin{figure}[H]
\centering
\begin{tikzpicture}[
    node distance=2.2cm,
    box/.style={rectangle, draw, text width=3cm, text centered, minimum height=1cm, fill=white},
    arrow/.style={->, >=stealth, thick}
]
\node[box] (contrib) {贡献者 \\ Contributor};
\node[box, right=of contrib] (declare) {公开声明 $g=0$ \\ Public $g=0$};
\node[box, right=of declare] (audit) {接受 $M>1$ 审计 \\ $M>1$ Audit};
\node[box, below=of audit] (vote) {内核投票 \\ Kernel Vote};
\node[box, left=of vote] (rotate) {进入轮换池 \\ Enter Rotation};

\draw[arrow] (contrib) -- (declare);
\draw[arrow] (declare) -- (audit);
\draw[arrow] (audit) -- (vote);
\draw[arrow] (vote) -- (rotate);
\draw[arrow, dashed] (rotate.south) -- ++(0,-1) -| (contrib.south);
\end{tikzpicture}
\caption{内核晋升五阶段路径 \\ Five-Stage Kernel Promotion Path}
\end{figure}

\subsection{阶段一：累积有效贡献}

候选人必须在贡献者层保持活跃，并至少累积一次被内核接受的有效贡献。
``有效贡献''的定义包括但不限于：
\begin{itemize}[nosep]
    \item 代码提交被合并至SCX主仓库
    \item 一篇 hostile review 论文被内核认可
    \item 发现并报告一个经验证的漏洞
    \item 成功复现一个已发布的审计结果
\end{itemize}

\subsection{阶段二：公开 $g=0$ 声明}

候选人必须提交并公开发布一份符合规范的$g=0$声明（详见第\ref{sec:g0-protocol}节）。
该声明是\textit{不可撤销的}——一旦签署并发布，候选人即永久性地放弃了以所有权为基础的
任何协议控制主张。

\subsection{阶段三：接受 $M>1$ 独立审计}

候选人的偏差参数$g$必须由至少$M > 1$名独立审计者进行评估。
审计者必须包括：
\begin{itemize}[nosep]
    \item 至少一名当前内核成员
    \item 至少一名非内核的贡献者
    \item 审计日志完全公开，任何人可复现
\end{itemize}

\begin{lemma}[$M>1$审计的检测下界]
设审计者数量为$M$，每位审计者对偏差$g \neq 0$的独立检测概率为$p_{\text{detect}}$。
则至少一名审计者检测到偏差的概率下界为：
\begin{equation}
    \Pbb(\text{至少一次检测} \mid g \neq 0) \geq 1 - (1 - p_{\text{detect}})^M
\end{equation}
当$M \geq 3$且$p_{\text{detect}} \geq 0.5$时，该概率超过$0.875$。
\end{lemma}

\subsection{阶段四：内核一致性投票}

内核对新成员加入实行\textbf{全票通过规则（Unanimous Consent Rule）}。
\textit{任何一名}现有内核成员拥有否决权——一票反对即阻止加入。
这一规则的设计理由在第\ref{sec:game-theory}节中的博弈论分析中给出。

\begin{definition}[内核投票函数]
设当前内核为$\kernelSet = \{k_1, \ldots, k_{\numMaintainers}\}$，候选人$c$的投票结果为：
\begin{equation}
    V(c) = \prod_{i=1}^{\numMaintainers} v_i(c), \quad v_i(c) \in \{0, 1\}
\end{equation}
其中$v_i(c) = 1$表示成员$k_i$同意，$V(c) = 1$当且仅当所有成员同意。
\end{definition}

\subsection{阶段五：进入轮换池}

通过投票后，候选人进入内核\textit{轮换池}。轮换池中的成员按照预定轮换日程
依次进入内核任职。任期结束后，成员自动退出内核并返回贡献者层——
但可以重新申请。这一机制确保内核始终保持新鲜血液。

\begin{govbox}
\textbf{为什么不是终身制？} 终身维护者是协议的癌症。任何人在长期掌控校准锚点后，
即使初始$g=0$，也会逐渐积累隐性偏差——不是因为恶意，而是因为人类认知的固有局限。
轮换机制不是对维护者的不信任——轮换是对\textit{人类生物学的尊重}。
\end{govbox}

% ===================================================================
% SECTION 4: THE g=0 DECLARATION PROTOCOL
% ===================================================================
\section{$g=0$宣言协议：所有权的数学放弃}
\section{The $g=0$ Declaration Protocol: Mathematical Renunciation of Ownership}
\label{sec:g0-protocol}

\subsection{为什么需要 $g=0$ 声明}

SCX协议的核心操作是\textit{审计}——对人工智能系统、科学声明、治理决策进行客观的、
可复现的偏差检测。如果审计者本身存在未被声明的偏差，整个审计框架的数学基础就会崩塌。
这是SCX的根本性自指困境：\textit{谁来审计审计者？}

SCX的回答是：审计者通过公开的$g=0$声明，将自己置于与被审计对象\textit{相同的审计框架之下}。
这不是道德承诺——这是数学约束。$g=0$声明使得审计者的行为也服从Hoeffding检测界，
从而使得审计者与被审计者处于博弈论上的对等地位。

\subsection{宣言的五项必要内容}

根据SCX社区宪法，一份有效的$g=0$声明必须包含以下五项内容。
缺少任何一项即构成无效声明，候选人不得进入内核。

\begin{enumerate}[leftmargin=*, label=\textbf{(\arabic*)}]
    \item \textbf{身份声明 (Identity Declaration)}: \\
    实名或可验证的长期化名。化名必须满足``长期可追溯''条件——即该化名在过去至少
    12个月内持续用于SCX相关的公开活动（代码提交、论文发表、社区讨论），
    且可被加密签名验证。

    \item \textbf{所有权放弃声明 (Ownership Renunciation)}: \\
    明确放弃对SCX协议的任何所有权主张。包括但不限于：
    \begin{itemize}[nosep]
        \item 不主张对SCX商标、代码、协议规范的所有权
        \item 不主张对CEC校准锚点值的排他性决定权
        \item 不主张对协议演化方向的个人否决权
    \end{itemize}

    \item \textbf{审计同意声明 (Audit Consent)}: \\
    同意接受定期的$M>1$独立偏差审计，且审计日志\textit{完全公开}。
    该同意不可撤回——即使声明者退出内核，其在任期间的所有行为仍然接受公开审计。

    \item \textbf{轮换同意声明 (Rotation Consent)}: \\
    同意轮换机制——任期$T$结束后自动退出内核。可以重新申请，
    但重新申请需经过完整的五阶段晋升流程（包括新的$M>1$审计）。

    \item \textbf{加密签名 (Cryptographic Signature)}: \\
    声明文件须使用声明者的GPG密钥签名，并发布于SCX公开Git仓库的
    \texttt{declarations/}目录下。签名确保声明的不可否认性和完整性。
\end{enumerate}

\begin{principle}
\textbf{$g=0$声明的不可逆原则。} 一旦签署并发布，$g=0$声明的五项内容不可单方面修改或撤回。
声明者可以退出社区，但退出行为不改变声明的历史约束力——
其在任期间的所有行为仍然受到声明中审计同意的约束。
\end{principle}

\subsection{$g=0$的博弈论解释}

$g=0$声明不仅仅是一份文件——它是一个博弈论信号。在Spence信号传递博弈的意义上，
$g=0$声明的\textit{成本}——永久放弃所有权、接受终身可审计性、同意轮换——
创造了分离均衡的条件：只有真正$g=0$的个体才愿意承担这一成本。

\begin{proposition}[$g=0$声明的分离均衡]
在审计博弈中，设候选人的真实偏差为$g_{\text{true}} \in \{0, \delta\}$，
其中$\delta > 0$。签署$g=0$声明的成本为$\kappa$（包括所有权放弃的永久成本和
审计暴露的风险成本）。存在唯一的分离均衡当且仅当：
\begin{equation}
    \kappa > \text{预期捕获收益} \times \Pbb(\text{捕获成功} \mid g = \delta)
\end{equation}
即：对于$g = \delta$的候选人，签署声明的期望成本超过期望收益，因此他们不会签署；
而$g = 0$的候选人因无隐藏偏差，$\kappa \approx 0$（真正的协议服务本身就是他们的收益），
因此他们愿意签署。
\end{proposition}

% ===================================================================
% SECTION 5: MAINTAINER ROTATION GAME THEORY
% ===================================================================
\section{维护者轮换博弈论：为什么定理审计的是维护者}
\section{Maintainer Rotation Game Theory: Why the Theorem Audits the Maintainer}
\label{sec:game-theory}

\subsection{轮换博弈的基本设定}

考虑一个由$\numMaintainers$名维护者组成的无限期重复博弈。在每个轮换周期
$T$内，每位维护者$i$选择其偏差参数$g_i \in [0, G_{\max}]$。
维护者的阶段收益为：
\begin{equation}
    \payoff{i}(g_i, \mathbf{g}_{-i}) = B(g_i) - \auditCost \cdot \Pbb(\text{检测到} \mid g_i) - P(g_i, T)
    \label{eq:payoff}
\end{equation}
其中：
\begin{itemize}[nosep]
    \item $B(g_i)$：偏差$g_i$带来的个人收益（如操纵CEC基准值以获取声誉或经济利益）
    \item $\auditCost$：被检测到偏差后的惩罚成本（包括强制退出内核和声誉损失）
    \item $\Pbb(\text{检测到} \mid g_i)$：在$M>1$互相审计下偏差被检测的概率
    \item $P(g_i, T)$：轮换结束时偏差暴露的概率——随着$T$增长而趋近于1
\end{itemize}

\subsection{Hoeffding界的轮换应用}

轮换机制的核心数学洞见是：\textit{有限任期使得隐藏偏差的概率随时间指数衰减}。
这是Hoeffding不等式的直接应用。

\begin{theorem}[轮换偏差检测定理]
设维护者$i$的真实偏差为$g_i$，每轮审计提供一次对$g_i$的独立估计
$\hat{g}_i^{(t)} = g_i + \varepsilon_t$，其中$\varepsilon_t$是零均值、
有界于$[-\Delta, \Delta]$的独立噪声。在$T$轮审计后，样本均值未检测到
$g_i > \Delta$的概率上界为：
\begin{equation}
    \Pbb\left(\left|\frac{1}{T}\sum_{t=1}^{T} \hat{g}_i^{(t)} - g_i\right| \leq \Delta \;\middle|\; g_i > 2\Delta\right)
    \leq 2e^{-2T\Delta^2 / (4\Delta^2)} = 2e^{-T/2}
    \label{eq:hoeffding}
\end{equation}
当$T \to \infty$时，$\Pbb(\text{未检测}) \to 0$。轮换通过限制$T$至有限值，
确保偏差必然在可计算的时间窗口内暴露。
\end{theorem}

\begin{proof}
由Hoeffding不等式，对于独立有界随机变量$X_1, \ldots, X_T$，其中
$X_t \in [a_t, b_t]$，有：
\begin{equation}
    \Pbb\left(\left|\frac{1}{T}\sum_{t=1}^{T}(X_t - \E[X_t])\right| \geq \varepsilon\right)
    \leq 2\exp\left(-\frac{2T^2\varepsilon^2}{\sum_{t=1}^{T}(b_t - a_t)^2}\right)
\end{equation}
令$X_t = \hat{g}_i^{(t)}$，$\E[X_t] = g_i$，$b_t - a_t \leq 2\Delta$，$\varepsilon = \Delta$，
代入即得式(\ref{eq:hoeffding})。
\end{proof}

\subsection{全票通过规则的博弈论基础}

内核采用全票通过规则（而非多数决）的深层原因在于：\textit{多数决可以被联盟操纵，
而全票通过规则使得任何一个诚实维护者可以单独阻止被捕获的联盟注入偏差}。

\begin{theorem}[全票通过的防捕获定理]
设内核$\kernelSet$中有$B$名已被捕获的维护者（$g > 0$），其试图引入第$B+1$名
被捕获者。在多数决规则下，所需$B \geq \lceil \numMaintainers/2 \rceil$。
在全票通过规则下，所需$B = \numMaintainers$。因此，只要$\numMaintainers - B \geq 1$
（即至少有一名诚实维护者），全票通过规则就阻止了捕获扩张。
\end{theorem}

\begin{hitbox}
\textbf{诚实暴击：为什么多数决在协议治理中是危险的。}
大多数人类组织使用多数决——不是因为多数决更好，而是因为多数决更快。
但协议治理不是民主选举。协议治理是\textit{审计完整性的维护}。
在审计完整性的问题上，没有``多数是对的而少数是错的''这回事——
\textit{一个诚实的反对票在数学上等价于一个被检测到的偏差}。
如果内核使用多数决，那么捕获51\%的内核成员就足以腐化整个协议。
全票通过规则将这一阈值提高到了100\%——使得捕获必须彻底完成才能成功，
而彻底捕获在公开审计和轮换机制下几乎不可能。
\end{hitbox}

\subsection{若干均衡分析}

\begin{proposition}[$M>1$互相审计的纳什均衡]
在$M>1$的互相审计博弈中，所有维护者报告$g_i = 0$且实际保持$g_i = 0$
构成一个纳什均衡。在此均衡下：
\begin{equation}
    \payoff{i}(0, \mathbf{0}_{-i}) = 0 \quad \text{（协议服务本身就是收益，$B(0) = 0$）}
\end{equation}
偏离至$g_i > 0$的单期收益为：
\begin{equation}
    \payoff{i}(g_i, \mathbf{0}_{-i}) = B(g_i) - \auditCost \cdot \Pbb(\text{检测} \mid g_i) - P(g_i, T)
\end{equation}
当$\auditCost$足够大且$T$足够小时，偏离收益为负，均衡成立。
\end{proposition}

\begin{proposition}[轮换博弈的子博弈完美性]
在无限期重复的维护者轮换博弈中，策略``在每轮任期内保持$g_i = 0$''
构成一个子博弈完美均衡——因为任何轮换任期都是有限期博弈，
而有限期博弈中偏离的唯一可能受益期为当前任期。由于轮换保证了
$\Pbb(\text{任期内被检测}) \xrightarrow{T \to \infty} 1$，
且退出后的审计同意（$g=0$声明第3项）消除了``最后一期偏离''的收益，
该策略在整个博弈树上是最优的。
\end{proposition}

% ===================================================================
% SECTION 6: DISCUSSION SPACES
% ===================================================================
\section{讨论空间：三层信息隔离架构}
\section{Discussion Spaces: Three-Layer Information Partitioning}

\subsection{信息分层原则}

不同治理层级需要不同的信息可见性，这不是秘密主义——这是\textit{安全工程}。
将内核讨论完全公开会导致两个问题：(1) 新维护者候选人的审计讨论包含敏感个人信息，
公开讨论可能违反隐私规范；(2) 安全漏洞的早期讨论如果完全公开，可能在修复发布前
被恶意利用。

\begin{table}[H]
\centering
\caption{三层讨论空间 \\ Three-Layer Discussion Spaces}
\renewcommand{\arraystretch}{1.3}
\begin{tabular}{p{3.5cm}p{3cm}p{5cm}}
\toprule
\textbf{层级} & \textbf{可见性} & \textbf{讨论内容} \\
\textbf{Layer} & \textbf{Visibility} & \textbf{Content} \\
\midrule
内核讨论 & 私有 & \makecell[l]{维护者轮换日程 \\ 新维护者审核 \\ CEC校准异常调查 \\ 审计标准修订提案} \\
Kernel & Private & \\
\midrule
贡献者讨论 & 半公开 & \makecell[l]{代码审查 \\ hostile review 论文 \\ 实验复现 \\ 新方向提案} \\
Contributors & Semi-public & \\
\midrule
公开讨论 & 完全公开 & \makecell[l]{使用问题 \\ \Cercis{} 分数比较 \\ 审计日志讨论 \\ 任何声称的独立验证} \\
Public & Fully open & \\
\bottomrule
\end{tabular}
\end{table}

\subsection{平台架构}

SCX社区使用去中心化、抗审查的工具栈，避免单一平台依赖性。

\begin{table}[H]
\centering
\caption{社区平台架构 \\ Community Platform Architecture}
\begin{tabular}{lll}
\toprule
\textbf{功能 Function} & \textbf{工具 Tool} & \textbf{选择理由 Rationale} \\
\midrule
代码 + 论文 & GitHub（现有仓库） & 版本控制、审计追踪 \\
内核私密讨论 & Signal 群组 & 端到端加密 \\
贡献者讨论 & Matrix & 开源、去中心化、联邦协议 \\
公开讨论 & GitHub Issues + Discussions & 公开、可搜索、可引用 \\
审计日志 & 公开 Git 仓库 & 不可篡改、可复现 \\
正式声明 & arXiv + GitHub Releases & 永久存档、版本化管理 \\
\bottomrule
\end{tabular}
\end{table}

\begin{govbox}
\textbf{为什么内核使用Signal而非Matrix？} Signal提供经过广泛安全审计的端到端加密。
内核讨论的内容（候选人审核、校准异常调查）如果泄露，可能被利用来破坏审计中立性。
当内核讨论涉及$M>1$审计结果时，对话内容的保密性与审计日志的公开性之间存在
\textit{时间差}——讨论是暂时的保密，审计日志是永久的公开。
讨论结束后，所有实质性决策必须以公开审计日志的形式发布。
\end{govbox}

% ===================================================================
% SECTION 7: CONFLICT RESOLUTION
% ===================================================================
\section{冲突解决协议：从僵局到裁决}
\section{Conflict Resolution Protocol: From Deadlock to Adjudication}

\subsection{全票通过制下的僵局问题}

全票通过规则（第\ref{sec:game-theory}节）具有强大的防捕获能力，但代价是
\textit{僵局风险}——当内核成员在某个议题上无法达成一致时，没有任何决策可以通过。
SCX社区宪法规定了三级升级路径来解决僵局。

\subsection{三级升级路径}

\begin{enumerate}[leftmargin=*, label=\textbf{Level \arabic*:}]
    \item \textbf{内核内部延长讨论 (Extended Kernel Deliberation)}: \\
    当内核在标准讨论期内无法达成一致时，自动触发72小时延长讨论期。
    在此期间，每位内核成员必须提交一份书面的反对理由或同意理由，
    发表于内核私有讨论组。延长讨论期不可跳过。

    \item \textbf{贡献者层公开讨论与非约束性投票 (Contributor-Layer Public Deliberation)}: \\
    若延长讨论期后仍然僵局，议题降级至贡献者层进行公开讨论。
    贡献者层讨论完全公开（即在GitHub Discussions中进行），
    所有贡献者可投票提供建议。该投票是\textit{非约束性}的——
    内核有权在考虑贡献者建议后自行决策，但必须在最终的审计日志中
    记录贡献者投票结果及其决策理由。

    \item \textbf{独立第三方 $M>1$ 审计 (Independent Third-Party $M>1$ Audit)}: \\
    若议题涉及$g=0$的判定（即争议焦点是某人是否$g \neq 0$），
    则要求被质疑方接受独立第三方的$M>1$偏差审计。
    审计者必须来自SCX社区外部（即非当前内核或贡献者成员），
    且审计结果完全公开。
\end{enumerate}

\begin{figure}[H]
\centering
\begin{tikzpicture}[
    node distance=2cm,
    decision/.style={diamond, draw, text width=2.5cm, text centered, aspect=2, fill=white},
    action/.style={rectangle, draw, text width=3cm, text centered, minimum height=1cm, fill=white},
    arrow/.style={->, >=stealth, thick}
]
\node[action] (deadlock) {僵局检测 \\ Deadlock};
\node[action, below=of deadlock] (level1) {Level 1: 延长讨论 \\ 72小时};
\node[decision, below=of level1] (resolved1) {解决?};
\node[action, right=of resolved1, xshift=1.5cm] (done1) {决议通过};
\node[action, below=of resolved1] (level2) {Level 2: 贡献者公开讨论};
\node[decision, below=of level2] (resolved2) {解决?};
\node[action, right=of resolved2, xshift=1.5cm] (done2) {决议通过};
\node[action, below=of resolved2] (level3) {Level 3: 独立$M>1$审计};

\draw[arrow] (deadlock) -- (level1);
\draw[arrow] (level1) -- (resolved1);
\draw[arrow] (resolved1) -- node[above] {是} (done1);
\draw[arrow] (resolved1) -- node[left] {否} (level2);
\draw[arrow] (level2) -- (resolved2);
\draw[arrow] (resolved2) -- node[above] {是} (done2);
\draw[arrow] (resolved2) -- node[left] {否} (level3);
\end{tikzpicture}
\caption{冲突解决三级升级路径 \\ Three-Level Conflict Resolution Escalation}
\end{figure}

\subsection{冲突解决的形式化属性}

\begin{definition}[冲突解决协议的活跃性]
冲突解决协议满足\textit{活跃性 (liveness)}当且仅当：对任何僵局议题，
存在一条有限路径（即Level 1、Level 2、或Level 3之一）最终产生确定性的决议。
\end{definition}

\begin{proposition}[冲突解决协议的活跃性证明]
三级升级路径满足活跃性——因为Level 3（独立第三方$M>1$审计）
总是终止（审计是有限时间操作），且产生一个确定性的$g$值估计。
该估计本身不强制决议——但它消除了信息不对称，使得剩余的内核成员
能够基于客观数据达成共识。
\end{proposition}

% ===================================================================
% SECTION 8: EMERGENCY PROCEDURES
% ===================================================================
\section{紧急程序：当守护者成为威胁}
\section{Emergency Procedures: When the Guardian Becomes the Threat}

\subsection{触发条件}

紧急程序是为应对内核成员$g \neq 0$或恶意行为而设计的快速响应机制。
正常的内核轮换和定期审计足以应对\textit{隐性偏差累积}——紧急程序针对的是
\textit{显性恶意行为}，包括但不限于：
\begin{itemize}[nosep]
    \item 未经共识的单方面修改CEC基准值
    \item 拒绝接受定期审计或拒绝公开审计日志
    \item 有证据表明篡改或伪造审计数据
    \item 利用内核地位谋取外部经济利益
\end{itemize}

\subsection{四阶段紧急响应协议}

\begin{enumerate}[leftmargin=*, label=\textbf{E\arabic*}]
    \item \textbf{触发 (Trigger)}: \\
    任何贡献者均可触发紧急审计。触发者必须在GitHub公开提交$g \neq 0$的
    初步证据（证据标准：合理怀疑，不要求完整证明）。触发一旦提交，
    即启动48小时倒计时。

    \item \textbf{48小时独立审计 (48-Hour Independent Audit)}: \\
    其余内核成员（不包括被质疑者）须在48小时内组织$M>1$独立审计。
    如果被质疑者是内核中的多数，可由贡献者层指定独立审计者。

    \item \textbf{裁决与执行 (Verdict and Execution)}: \\
    若审计确认$g \neq 0$：（a）被质疑者立即移出内核；
    （b）完整审计日志公开发布；（c）所有被质疑者经手的校准记录标记为待审查。
    若审计确认$g = 0$：触发者不会受到惩罚——善意的错误指控是审计生态的一部分。

    \item \textbf{继任 (Succession)}: \\
    从贡献者轮换池中提名替代者。如果轮换池中有已通过审计的候选人，
    可直接进入内核投票；否则，启动加速晋升流程。
\end{enumerate}

\begin{hitbox}
\textbf{诚实暴击：为什么善意误报不应当受到惩罚。}
如果触发紧急审计的人可能因为误报而受到惩罚，那么就没有人会触发紧急审计——
这正是恶意维护者所希望的。惩罚善意误报就是奖励恶意行为。
SCX社区宪法的设计原则是：宁可容忍一千次误报，也不让一次真正的恶意逃过审计。
这不是宽容——这是冷酷的博弈论计算。
\end{hitbox}

\subsection{紧急程序的时间约束}

\begin{theorem}[紧急程序的响应上界]
从触发到裁决的最长时间为72小时（触发后立即启动48小时审计窗口，
审计完成后24小时内发布裁决）。此时间约束确保了恶意维护者在
被检测后没有足够时间造成不可逆的损害。
\end{theorem}

% ===================================================================
% SECTION 9: FUNDING AND SUSTAINABILITY
% ===================================================================
\section{资金来源与可持续模型：不从社区收钱}
\section{Funding and Sustainability: No Money from the Community}

\subsection{核心原则：社区不是收入来源}

SCX社区宪法规定了一个根本性的经济约束：\textbf{社区不从社区成员收钱}。
这不是道德姿态——这是防止经济激励腐化审计中立性的结构保障。
如果社区的运营资金依赖于社区成员的缴费，那么社区管理层就有了
\textit{保留缴费者的隐性激励}——这会在审计中产生系统性偏差。

\begin{table}[H]
\centering
\caption{资金来源结构 \\ Funding Source Architecture}
\renewcommand{\arraystretch}{1.3}
\begin{tabular}{>{\bfseries}ll}
\toprule
\textbf{来源 Source} & \textbf{用途 Purpose} \\
\midrule
\Yajie{} API 校准费 & 内核维护者基本运维 \\
Yajie API Calibration Fees & Basic Kernel Maintainer Operations \\
\midrule
自愿捐赠 & 贡献者激励 \\
Voluntary Donations & Contributor Incentives \\
\midrule
学术基金 & hostile review 奖金 \\
Academic Grants & Hostile Review Bounties \\
\midrule
SCX Prize 奖池 & 年度审计贡献奖 \\
SCX Prize Pool & Annual Audit Contribution Award \\
\bottomrule
\end{tabular}
\end{table}

\subsection{资金来源的独立性分析}

\begin{proposition}[资金来源独立性条件]
社区的资金来源满足\textit{独立性条件}：没有任何单一资金来源占总收入超过50\%，
且没有任何资金来源对审计结果有直接利益关系。这确保了：
\begin{equation}
    \forall \text{ 资金来源 } s : \frac{\text{收入}_s}{\text{总收入}} < 0.5
    \;\wedge\; \text{利益冲突}_s = \emptyset
\end{equation}
\end{proposition}

\begin{govbox}
\textbf{\Yajie{} API校准费的经济逻辑。} \Yajie{} API的企业用户支付校准费来获取
准确的CEC基准值。这些费用流向维护者——不是作为``利润''，而是作为校准工作的
\textit{运营成本补偿}。维护者的经济利益与CEC基准值的\textit{准确性}绑定，
而不是与特定数值绑定。这是关键设计：如果维护者试图操纵基准值，他们会
\textit{损失}企业客户的信任和费用，而不是获得更多费用。这就是$\sum g = 0$的
经济基础。
\end{govbox}

% ===================================================================
% SECTION 10: AUDIT AND VERIFICATION FRAMEWORK
% ===================================================================
\section{审计与验证框架：社区的可审计性}
\section{Audit and Verification Framework: Community Auditability}

\subsection{社区自身的可审计性}

SCX社区的核心产品是\textit{审计}——因此社区自身必须是\textit{可审计的}。
这是一个自指约束：社区不能要求被审计对象接受的标准，高于社区自身愿意接受的标准。

\begin{definition}[社区可审计性]
SCX社区满足\textit{完全可审计性 (Full Auditability)}当且仅当：
\begin{enumerate}[nosep, label=(\roman*)]
    \item 所有内核决策都有对应的公开审计日志
    \item 所有维护者的$g$参数估计定期公开发布
    \item 所有校准锚点的修改历史完全公开且可追溯
    \item 任何第三方可独立复现所有审计结果
\end{enumerate}
\end{definition}

\subsection{审计日志规范}

所有审计日志须存储在公开Git仓库中，格式满足以下要求：

\begin{itemize}[nosep]
    \item \textbf{不可篡改性 (Immutability)}: 每次提交包含SHA-256哈希链，
      确保历史记录不可被单方面修改
    \item \textbf{可复现性 (Reproducibility)}: 包含完整的输入数据引用、
      评估参数和算法版本号
    \item \textbf{时间戳 (Timestamping)}: 每次审计包含可信时间戳（RFC 3161）
    \item \textbf{签名 (Signatures)}: 每条审计记录由审计者GPG签名
\end{itemize}

\begin{proposition}[审计日志的篡改检测]
设审计日志是哈希链$H_0, H_1, \ldots, H_n$，其中$H_i = \text{SHA-256}(H_{i-1} \| D_i)$，
$D_i$为第$i$条审计记录。篡改任何$D_k$会导致$H_k, H_{k+1}, \ldots, H_n$全部改变，
而诚实节点维护的$H_n$将揭示不一致性。篡改被检测的概率为1（在计算绑定假设下）。
\end{proposition}

% ===================================================================
% SECTION 11: FIRST STEPS
% ===================================================================
\section{第一步：从零到一启动社区}
\section{First Steps: Bootstrapping the Community from Zero to One}

\subsection{四阶段启动路线图}

社区的启动不是一蹴而就的——它是一个逐步建立信任、规范和制度的过程。
以下四阶段路线图定义了从当前状态到完全运行的SCX社区的路径。

\begin{enumerate}[leftmargin=*, label=\textbf{Phase \arabic*:}]

    \item \textbf{建立初始内核 (Establish Initial Kernel)}: \\
    创建Signal群组，邀请2--5名已通过$M>1$审计的初始维护者。
    初始维护者的选择标准：已对SCX协议有实质性贡献（代码、论文、审计），
    且已通过非正式的偏差评估。初始内核的第一个任务是制定正式的内核章程。

    \item \textbf{发布维护者候选人声明模板 (Publish Candidate Declaration Template)}: \\
    在SCX公开GitHub仓库中发布标准化的$g=0$声明模板和GPG签名指南。
    模板包含第\ref{sec:g0-protocol}节定义的五项必要内容，
    并附有填写示例和常见问题解答。任何人均可提交——不需要邀请。

    \item \textbf{首次公开 hostile review 活动 (First Public Hostile Review Event)}: \\
    悬赏寻找SCX论文中的漏洞。赏金由学术基金和SCX Prize奖池提供。
    目标：验证社区审计能力、吸引高质量的贡献者、
    并通过实战证明$g=0$审计文化的可行性。

    \item \textbf{建立轮换日程 (Establish Rotation Schedule)}: \\
    为初始维护者设定第一轮任期和审计时间表。初始任期建议为6个月
    （比正常任期短，以便快速验证轮换机制的可行性）。
    同时启动贡献者轮换池的创建。

\end{enumerate}

\begin{table}[H]
\centering
\caption{社区启动时间线 \\ Community Bootstrapping Timeline}
\begin{tabular}{ccl}
\toprule
\textbf{阶段} & \textbf{时间} & \textbf{可交付物} \\
\textbf{Phase} & \textbf{Timeline} & \textbf{Deliverable} \\
\midrule
1 & 第1周 & 初始内核Signal群组建立 \\
2 & 第1--2周 & $g=0$声明模板发布 \\
3 & 第3--8周 & 第一次hostile review活动 \\
4 & 第4--8周 & 轮换日程发布，贡献者轮换池开放 \\
\bottomrule
\end{tabular}
\end{table}

% ===================================================================
% SECTION 12: CORE PRINCIPLES
% ===================================================================
\section{核心原则：社区宪法的不可动摇支柱}
\section{Core Principles: The Immovable Pillars of the Community Constitution}

\subsection{十二条宪法原则}

以下十二条原则构成SCX社区宪法的不可修改条款。任何修改这些原则的提议
本身即构成$g \neq 0$的证据——因为此类提议必然服务于修改者的利益
而非协议的利益。

\begin{enumerate}[leftmargin=*, label=\textbf{原则 \arabic*:}]

    \item \textbf{定理主权原则 (Theorem Sovereignty)}: \\
    社区不是任何个人的——社区是定理的。初始维护者只是第一批守护者。
    定理不依赖任何个人。任何时候社区可以审计维护者、替换维护者。

    \item \textbf{$g=0$至上原则 ($g=0$ Supremacy)}: \\
    所有维护者必须在数学上可证明的意义上保持$g=0$。$g=0$不是道德声明——
    它是审计博弈的纳什均衡条件。任何$g \neq 0$的行为，无论动机多么崇高，
    都在数学上破坏了协议的中立性。

    \item \textbf{公开审计原则 (Public Audit)}: \\
    所有审计日志完全公开。秘密审计不是审计——
    秘密审计是无法被第三者复现的独白。

    \item \textbf{轮换不可协商原则 (Rotation Non-Negotiability)}: \\
    维护者是轮换的。没有终身维护者。这不是对人的不信任——
    这是对生物学和博弈论的尊重。

    \item \textbf{全票通过原则 (Unanimous Consent)}: \\
    内核决策采用全票通过制。一个诚实的反对票在数学上等价于一个被检测到的偏差。

    \item \textbf{零所有权原则 (Zero Ownership)}: \\
    没有人拥有SCX协议。社区是一个维护网络，不是一个产权实体。
    任何对协议的所有权主张自动使其主张者丧失维护者资格。

    \item \textbf{社区自审计原则 (Community Self-Audit)}: \\
    社区自身接受与被审计对象相同的审计标准。审计者必须可被审计。

    \item \textbf{非盈利原则 (Non-Profit)}: \\
    社区不从社区成员收钱。经济利益只能来自外部（API费用、学术基金、捐赠），
    且来源多元化以确保独立性。

    \item \textbf{开放性不可逆原则 (Irreversible Openness)}: \\
    观察者层级永不对任何人关闭。零门槛进入是社区生命力的制度保障。

    \item \textbf{紧急审计不可惩罚原则 (Emergency Audit Non-Retaliation)}: \\
    触发紧急审计的人，即使最终审计结果为$g=0$，也不受惩罚。
    惩罚善意误报就是奖励恶意行为。

    \item \textbf{审计日志哈希链原则 (Audit Log Hash Chain)}: \\
    所有审计日志以SHA-256哈希链形式存储。任何篡改必然被检测。

    \item \textbf{社区高于维护者原则 (Community over Maintainer)}: \\
    如果必须在保护社区和保护特定维护者之间选择——选择保护社区。
    社区可以失去任何个人而存活。没有人如此重要以至于其利益高于定理的完整性。

\end{enumerate}

\begin{principle}
\textbf{元原则：宪法不可修改的本质条款。} 以上十二条原则构成社区宪法的
\textit{永恒条款 (Eternal Clauses)}。任何修改这些条款的提议不是``修宪''——
它是\textit{对宪法的攻击}，触发紧急审计程序。协议的治理规则可以演化——
但治理规则的\textit{目的}（保护定理的中立性、可审计性和可复现性）是永恒的。
\end{principle}

% ===================================================================
% SECTION 13: THE COMMUNITY IS THEOREM'S
% ===================================================================
\section{社区属于定理：最终声明}
\section{The Community Belongs to the Theorem: Final Declaration}

\subsection{超越个人的协议}

SCX协议的核心创新不是任何特定的算法、定理或工具。核心创新是\textit{一个约束}：
任何声称具有权威性的声明——无论是科学结论、AI系统输出还是治理决策——
都必须接受多专家、可复现、公开的因果偏差审计。

这个约束同样适用于SCX社区自身。社区的权威性必须可审计。
社区的结构必须阻止任何个人的偏差累积。社区的经济必须独立于社区的成员。

\subsection{守护者的悖论与解决}

\begin{hitbox}
\textbf{诚实暴击：守护者悖论。}
谁守护守护者？\textit{(Quis custodiet ipsos custodes?)} 罗马诗人尤维纳利斯
在两千年前提出了这个问题。SCX协议的回答是：\textbf{公开的、可复现的数学}。
守护者被定理守护——定理审计守护者，正如定理审计一切。维护者可以替换，
审计可以复现，轮换可以强制执行。不需要信任任何人——只需要相信
$e^{-2M\Delta^2}$ 在$M$足够大时趋近于零。
\end{hitbox}

\begin{govbox}
\textbf{治理透视：为什么这很重要。}
人类组织的历史是一部捕获与腐化的历史。公司被股东捕获，政府被利益集团捕获，
学术期刊被引用指标捕获，开源项目被终身仁慈独裁者捕获——然后独裁者离开或腐化。
SCX社区宪法的设计目标是打破这个循环。不是通过找到更好的独裁者——
而是通过\textit{消除独裁者这个职位}。定理是唯一的权威。维护者只是定理的仆人。
\end{govbox}

\subsection{社区宪法的签署}

本宪法文件本身是SCX社区的基础法律文件。它不是一个描述性文档——
它是一个\textit{规范性}文档。社区按照本文定义的规则运行，或者根本不能称为
SCX社区。

\begin{center}
\fbox{\begin{minipage}{0.85\textwidth}
\centering
\color{scxblue}
\textbf{SCX协议社区宪法}

\vspace{6pt}
{\small 本宪法于2026年7月2日由SCX协议社区治理工作组起草。\\
宪法生效需要：\\
(1) 至少2名初始内核成员签署本文件的GPG签名版本；\\
(2) 签名版本发布在SCX公开Git仓库的 \texttt{constitution/} 目录下。\\
宪法修改需要：全体现任内核成员的一致同意，\\
且修改不得违反第\ref{sec:core-principles}节定义的十二条宪法原则。}

\vspace{6pt}
{\small SCX Protocol Community Constitution \\
Drafted by SCX Protocol Community Governance Working Group, July 2, 2026. \\
Effective upon: (1) GPG-signed ratification by at least 2 initial kernel members; \\
(2) Publication in the SCX public Git repository under \texttt{constitution/}. \\
Amendments require: unanimous consent of all current kernel members, \\
and must not violate the twelve constitutional principles defined in Section~\ref{sec:core-principles}.}
\end{minipage}}
\end{center}

\vspace{12pt}

\begin{flushright}
\textit{——社区属于定理。定理不依赖任何个人。} \\
\textit{--- The community belongs to the theorem. The theorem depends on no individual.}

\vspace{24pt}
SCX Protocol Community Governance Working Group \\
2026年7月2日 / July 2, 2026
\end{flushright}

% ===================================================================
% REFERENCES
% ===================================================================
\newpage
\begin{thebibliography}{99}

\bibitem{scx_theory}
SCX. \textit{SCX: Structured Causal eXamination — A Mathematical Framework for Multi-Expert Causal Audit}. SCX Internal Paper, 2026.

\bibitem{scx_protocol_gov}
SCX Protocol Governance Research Division. \textit{SCX Protocol Governance: Maintainer Rotation Game Theory and the Trustless Architecture}. SCX Internal Paper, 2026.

\bibitem{scx_maintainer_analysis}
SCX Protocol Governance Research Division. \textit{SCX Protocol Maintainer Candidate Analysis: Pre-Audit Assessment of Four Candidates for $g=0$}. SCX Internal Paper, 2026.

\bibitem{yajie}
SCX. \textit{Yajie Protocol: Multi-Expert Nash-Pareto Equilibrium for Causal Bias Audit}. SCX Internal Paper, 2026.

\bibitem{spring}
SCX. \textit{Spring Configuration Space Convergence: A Fixed-Point Theory of Iterative Representation Refinement}. SCX Internal Paper, 2026.

\bibitem{hoeffding}
W.~Hoeffding. ``Probability Inequalities for Sums of Bounded Random Variables.'' \textit{Journal of the American Statistical Association}, 58(301):13--30, 1963.

\bibitem{spence}
M.~Spence. ``Job Market Signaling.'' \textit{The Quarterly Journal of Economics}, 87(3):355--374, 1973.

\bibitem{fudenberg_tirole}
D.~Fudenberg and J.~Tirole. \textit{Game Theory}. MIT Press, 1991.

\bibitem{ostrom}
E.~Ostrom. \textit{Governing the Commons: The Evolution of Institutions for Collective Action}. Cambridge University Press, 1990.

\bibitem{lessig}
L.~Lessig. ``Code Is Law: On Liberty in Cyberspace.'' \textit{Harvard Magazine}, 2000.

\bibitem{nakamoto}
S.~Nakamoto. ``Bitcoin: A Peer-to-Peer Electronic Cash System.'' 2008.

\bibitem{buterin}
V.~Buterin. ``Decentralized Autonomous Organizations and Governance.'' Ethereum Blog, 2014.

\end{thebibliography}

\end{document}
